\section{Platform context and limits of comparison}
\label{sec:platforms}
Published platform results provide reference points for the next evaluation, not a ranking against this study. Table~\ref{tab:platforms} records vendor-reported LongMemEval headline scores inspected on 22 October 2026 (UTC). We did not rerun these systems. Their source scope, readers, inference budgets, and grading differ from the sixty-case Cortex experiment.

\begin{table}[htbp]
\centering\small
\begin{tabularx}{\linewidth}{@{}L{.16\linewidth}rY@{}}
\toprule
Platform & Reported (\%) & Reported scope and protocol qualification\\
\midrule
Mem0 & 94.4 & Managed platform, top-200 retrieval; mean 6,787 context tokens. Proprietary optimizations are excluded from the open-source SDK \citep{mem0eval}.\\
Zep & 90.2 & 451/500 questions; multiple retrieval scopes and a cross-encoder reranker. Reader and judge are GPT-5.4; reader uses medium reasoning \citep{zepeval}.\\
Hindsight & 94.6 & LongMemEval-S headline on the project overview. Its recall and agentic reflection paths have different call budgets; the headline alone does not identify a matched configuration \citep{hindsightdocs}.\\
\bottomrule
\end{tabularx}
\caption{External, vendor-reported results, unverified by this study. Cortex's 78.3\% reference alignment includes conditional answers on a different sixty-case sample and is deliberately excluded from this table. No percentage-point gap or superiority claim can be inferred.}
\label{tab:platforms}
\end{table}

Mem0 describes additive temporal metadata and time-sensitive retrieval ranking \citep{mem0temporal}. Zep's Graphiti represents changing entity relationships in a temporal knowledge graph \citep{graphitidocs}. Hindsight describes semantic, keyword, graph, and temporal retrieval, followed by reranking, as well as a separate reflection operation \citep{hindsightdocs}. These are concrete design comparisons. Their benefit for a particular workload still requires an experiment.

Cortex emphasizes a small Go/PostgreSQL package with explicit source support and state transitions. It has lexical and semantic retrieval, contradiction history, activation decay, and source-linked date arithmetic. It does not yet have a validated learned retention policy, a general event-identity model, or an evaluated graph-based retrieval stage. Its engineering checks establish observed behavior under tests; the current evidence does not establish competitive answer quality, lower cost, or lower latency than the platforms above.
